%% plantilla.tex — Presentacion de ejemplo para Simulacion de Sistemas (ITBA)
%%
%% Estructura obligatoria segun la Guia de la catedra (../AGENTS.md §5.3 / docs):
%%   1. Introduccion / Sistema Real / Fundamentos  (< 1 min, max 3 diapositivas)
%%   2. Implementacion  (solo motor de simulacion; sin post-proceso ni I/O)
%%   3. Simulaciones    (sistema particular, observables definidos, repeticiones/tiempos)
%%   4. Resultados      (animacion -> evolucion temporal -> input vs observable, por parametro)
%%   5. Conclusiones    (1 diapositiva, solo sobre lo mostrado)
%%   +  Cierre "Muchas Gracias"
%%
%% Reglas de figuras: sin titulo interno ni caption; parametros fijos al costado
%% (usar \figuraparam). Ejes en palabras + unidades; leyendas SIN titulo (el modelo
%% y los parametros van al costado, no dentro del grafico).
%%
%% Compilar:  pdflatex plantilla.tex   (dos veces)

\documentclass[aspectratio=169]{beamer}
\usepackage{sdsbeamer}

% --- Datos de la presentacion (completar) ---
\title{Titulo del Trabajo Practico}
\subtitle{Simulacion de Sistemas --- ITBA}
\author[Grupo 2]{Arancibia Nicolás -- Legajo pendiente \and
  Morales Agustín -- Legajo pendiente \and Ronda Agustín -- Legajo 64507}
\institute[ITBA]{Grupo 2 --- Comisión S2 --- Instituto Tecnologico de Buenos Aires}
\date{\today}

\begin{document}

\begin{frame}[plain]
  \titlepage
\end{frame}

% ==========================================================================
\section{Introduccion}
% Max 3 diapositivas: sistema real + ecuaciones del modelo matematico general.
% ==========================================================================
\begin{frame}{Sistema real}
  \begin{itemize}
    \item Descripcion somera del sistema real a simular.
    \item Fenomeno de interes y por que amerita una simulacion.
  \end{itemize}
\end{frame}

\begin{frame}{Modelo matematico}
  Ecuaciones generales del modelo (sin sistema particular ni metodo de resolucion):
  \begin{equation}
    \vect{r}_i(t+\Delta t) = \vect{r}_i(t) + \vect{v}_i(t)\,\Delta t
  \end{equation}
  donde $\vect{r}_i$ es la posicion de la particula $i$ y $\vect{v}_i$ su velocidad.
  \begin{itemize}
    \item Vectores en negrita recta ($\vect{v}$); escalares en italica ($t$).
    \item Unidades MKS y notacion cientifica: $\Delta t = \sci{1.0}{-2}\uni{s}$.
  \end{itemize}
\end{frame}

% ==========================================================================
\section{Implementacion}
% Solo el motor de simulacion: arquitectura, pseudocodigo, UML. Sin post-proceso.
% ==========================================================================
\begin{frame}{Motor de simulacion}
  \begin{itemize}
    \item Arquitectura del motor (modulos, estructuras de datos).
    \item Pseudocodigo o diagrama del paso de simulacion.
    \item Sin post-proceso ni formatos de archivo I/O (§2.2 de la guia).
  \end{itemize}
\end{frame}

% ==========================================================================
\section{Simulaciones}
% Sistema particular, observables definidos matematicamente, repeticiones y tiempos.
% ==========================================================================
\begin{frame}{Sistema particular y observables}
  \begin{columns}[t]
    \begin{column}{0.5\textwidth}
      \textbf{Configuracion}
      \begin{itemize}
        \item Geometria y condiciones de contorno.
        \item Parametros fijos.
        \item Input variable; output a estudiar.
      \end{itemize}
    \end{column}
    \begin{column}{0.5\textwidth}
      \textbf{Observable} (definicion explicita):
      \begin{equation}
        \langle O \rangle = \frac{1}{N}\sum_{i=1}^{N} o_i
      \end{equation}
      Repeticiones: 50 seeds distintas. \\
      Tiempo simulado: $T = 2000$ pasos.
    \end{column}
  \end{columns}
\end{frame}

% ==========================================================================
\section{Resultados}
% Por cada parametro: animacion -> evolucion temporal -> input vs observable.
% ==========================================================================
\begin{frame}{Animacion caracteristica}
  % En el PDF entregado: fotograma fijo + link (no video embebido, §2.4/5.4).
  \figuraparam{example-image}{%
    \textbf{Parametros fijos}
    \begin{itemize}
      \item $L = 10\uni{m}$
      \item $\rho = 2$
      \item $\eta = 0.5\uni{rad}$
    \end{itemize}
    \vspace{0.5em}
    Video: \url{https://youtu.be/XXXXXXX}
  }
\end{frame}

\begin{frame}{Evolucion temporal del observable}
  % Solo evoluciones tipicas (valores extremos) para validar la metrica.
  % Leyenda SIN titulo: solo identifica cada curva; el modelo/rho van al costado.
  \figuraparam{example-image}{%
    \textbf{Vicsek}, $\rho = 2$ \\[0.4em]
    Marca vertical: inicio del estado estacionario. \\[0.4em]
    Escalar caracteristico: promedio temporal en el estacionario.
  }
\end{frame}

\begin{frame}{Input vs observable}
  % Promedios con barras de error; puntos con simbolos; rectas solo como guia.
  \figuraparam{example-image}{%
    \textbf{Vicsek} \\[0.4em]
    Promedio sobre 50 seeds. \\
    Barras de error: desvio estandar muestral. \\[0.4em]
    Una curva por densidad $\rho$.
  }
\end{frame}

% ==========================================================================
\section{Conclusiones}
% Solo sobre resultados mostrados; sin hipotesis no probadas ni pendientes.
% ==========================================================================
\begin{frame}{Conclusiones}
  \begin{itemize}
    \item Conclusion 1, basada en los resultados mostrados.
    \item Conclusion 2, basada en los resultados mostrados.
  \end{itemize}
\end{frame}

% ==========================================================================
\begin{frame}[plain]
  \centering
  \vfill
  {\Huge Muchas gracias}
  \vfill
\end{frame}

\end{document}
